\subsection{Part F(c). Two punctured elliptic curves, with exact local dimensions}

This subsection implements the successive examination of splittings proposed
in the closing pages of \S51, folio 734, scan 173, and folio 736, scan 178.
It concerns the geometric group of type $(1,1)$ occurring in \S52,
folios 743--746, scans 191--197. It supplies calculations for the abelian,
nilpotent and free metabelian unipotent approximations, with their Galois
actions. Splittings always mean conjugacy classes of splittings satisfying
the specified local conditions. Their parameter schemes can have positive
dimension even when the set of integral points is finite.

Let $E_d:y^2=x^3-d^2x$, $X_d=E_d\setminus\{O\}$, and take $d=1,6$.
We use $p=5$ and the rational base point $b=(0,0)$. Let $\mathcal E_d$
be the minimal proper regular model over $\mathbb Z$ and let
$\mathcal X_d=\mathcal E_d\setminus\overline{\{O\}}$. This model contains
the integral affine points on the displayed Weierstrass equation. We allow
a finite set $S$ of denominator primes, with $5\notin S$, and put $s=|S|$.
The bad primes are $B_1=\{2\}$ and $B_6=\{2,3\}$; put
$T=S\cup B_d\cup\{5\}$. They are not automatically counted in $s$:
at a bad prime outside $S$ the integral local condition is retained.

Write $U_n=\pi_1^{\rm un}(X_{d,\overline{\mathbb Q}},b)/\gamma_{n+1}$,
over $\mathbb Q_5$, and $V=T_5E_d\otimes\mathbb Q_5$. Thus $U_1=V$.
Write $M_n$ for the corresponding class-$n$ quotient of the free metabelian
quotient $U/[U',U']$. Both towers have $L_1=V$, $L_2=\mathbb Q_5(1)$.
For $U$ the higher central layers are $L_n=\operatorname{Lie}_n(V)$;
for $M$ write $L_n^{\rm met}$.
For each finite reduction type $\Sigma$, impose at $\ell\ne5$ the
Zariski closure of the local Kummer image of the indicated integral
component, or of the indicated punctured component at $\ell\in S$.
At $5$ impose crystallinity. Denote the resulting finite union of
refined Selmer schemes by $\operatorname{Sel}_{S,n}$, or
$\operatorname{Sel}^{\rm met}_{S,n}$. Its dimension means the maximum
dimension of its nonempty components. This convention permits the cuspidal
local component at every prime in $S$, even when it has no global integral
point. These schemes still contain the Kummer images of all
$\mathcal X_d(\mathbb Z[1/S])$.

\begin{statement}{F.6}{Proved}{The examples have ranks zero and one, and the usual Selmer ranks agree}
The curves $E_1$ and $E_6$ have Mordell--Weil ranks $0$ and $1$, respectively.
They both have good ordinary reduction at $5$, with $a_5=-2$.
Their Tate--Shafarevich groups are finite; consequently
\[
 \dim H_f^1(\mathbb Q,V)=r,\qquad r=0\ \hbox{or}\ 1,
 \qquad \dim H_f^1(\mathbb Q_5,V)=1.
\]
Here $H_f^1$ uses the Bloch--Kato local conditions at every place.
\end{statement}
\begin{proof}
The pure $2$-Selmer rank of a squarefree congruent-number curve is the
nullity of its Monsky matrix. In the additive Legendre-symbol convention,
for $d=2^\epsilon P_1\cdots P_t$ the even matrix is
\[
 M_e=\begin{pmatrix}D_2&A+D_2\\ A^{\mathsf T}+D_2&D_{-1}\end{pmatrix},
 \qquad
 A_{ij}=\left[\frac{P_j}{P_i}\right]\ (i\ne j),\quad
 A_{ii}=\sum_{j\ne i}A_{ij},
 \quad (D_h)_{ii}=\left[\frac h{P_i}\right].
\]
The empty odd matrix for $d=1$ has nullity zero. For $d=6$, $t=1$,
$P_1=3$, $A=0$, $D_2=D_{-1}=(1)$, so
$M_e=\left(\begin{smallmatrix}1&1\\1&1\end{smallmatrix}\right)$ has
nullity one. Thus the ranks are at most zero and one. This is the
explicit descent formula in \citet[\S2, (2.1)--(2.3)]{FE_DasSaikia2024},
attributed there to Monsky's appendix to Heath--Brown.

For $E_6$, $P=(12,36)$ satisfies
\[
 2P=(25/4,-35/8),\qquad
 4P=(1442401/19600,1726556399/2744000),\qquad 8P\ne O.
\]
Direct counts give $|E_6(\mathbb F_5)|=|E_6(\mathbb F_7)|=8$.
The kernel of reduction at a good prime is a pro-prime formal group, so
reduction is injective on torsion of order prime to that prime. Combining
the two primes shows that the order of rational torsion divides $8$.
Thus $P$ is not torsion, and the rank is one. The same counts at $3,5$
for $E_1$ are $4,8$; its four visible rational $2$-torsion points therefore
exhaust $E_1(\mathbb Q)$. The counts at $5$ also give $a_5=6-8=-2$ for
both curves, proving ordinarity.

The $2$-descent alone does not prove finiteness of the $5$-primary
Tate--Shafarevich group. We use a separate verified input. The primary
database entries $32.a3$ and $576.c3$ record analytic ranks zero and one
\citep{FE_LMFDBE1,FE_LMFDBE6}. Their conductors are $32$ and $576$.
The database's reliability statement specifies that analytic ranks at
most three at conductor at most $500000$ are rigorous modular-symbol
computations \citep{FE_LMFDBReliability}. Kolyvagin's own
\citet[Theorem 1, printed p.~4; also Example 1, p.~11]{FE_Kolyvagin1990}
states that, for a modular elliptic curve of analytic rank at most one,
the algebraic rank agrees and the Tate--Shafarevich group is finite.
The hypotheses hold here. The Kummer exact sequence after tensoring
with $\mathbb Q_5$ gives the global assertion; the formal logarithm gives
the one-dimensional local assertion.

The program \texttt{code/f\_elliptic\_dimensions.py} checks both descent
matrices, all displayed point operations, and all four reduction counts
by exact rational or finite-field arithmetic. The analytic-rank input is
identified separately in its output and is not purported to be recomputed.
\end{proof}

\begin{statement}{F.7}{Proved}{Full and metabelian local dimensions through class eight}
For either curve the exact local dimensions are
\[
 \begin{array}{c|rrrrrrrr}
 n&1&2&3&4&5&6&7&8\\\hline
 \dim H_f^1(\mathbb Q_5,U_n)&1&2&4&7&13&22&40&70\\
 \dim H_f^1(\mathbb Q_5,M_n)&1&2&4&7&11&16&22&29
 \end{array}
\]
The complete linear data used below are
\[
 \begin{array}{c|rrrrrrrr}
 n&1&2&3&4&5&6&7&8\\\hline
 w_n=\dim L_n&2&1&2&3&6&9&18&30\\
 \operatorname{tr}(c\mid L_n)&0&-1&0&-1&0&-1&0&-2\\
 \dim L_n^{c=-1}&1&1&1&2&3&5&9&16\\
 \dim L_n^{\rm met}&2&1&2&3&4&5&6&7\\
 \dim (L_n^{\rm met})^{c=-1}&1&1&1&2&2&3&3&4
 \end{array}.
\]
The degree is the lower-central degree, equal here to the negative
geometric weight. It is not the doubled Tate grading of the tripod.
\end{statement}
\begin{proof}
The topological fundamental group of a once-punctured torus is free on
two generators. The central graded Lie algebra of its unipotent
completion is therefore the free Lie algebra on $V$. The Hodge
filtration on $V$ has one generator in degree $0$ and one in degree
$-1$. A Lie word of length at least two in just the degree-zero
generator vanishes. Hence $F^0L_n=0$ for $n\ge2$, whereas
$\dim F^0V=1$. The de Rham comparison therefore gives
\[
 \dim(U_n^{\rm dR}/F^0)=\sum_{m=1}^n w_m-1,
 \qquad w_m=\frac1m\sum_{e\mid m}\mu(e)2^{m/e}.
\]
There are no Frobenius-fixed vectors in the pure negative-weight
layers. The crystalline torsor classification identifies the displayed
de Rham quotient with $H_f^1(\mathbb Q_5,U_n)$; alternatively its
dimension follows by successive central extensions and
\citet[Corollary 3.1.11]{FE_Betts2023}.

Complex conjugation on $V$ has eigenvalues $1,-1$. The character form
of the Witt formula is
\[
 \operatorname{tr}(c\mid L_m)
 =\frac1m\sum_{e\mid m}\mu(e)
       \bigl(\operatorname{tr}(c^e\mid V)\bigr)^{m/e},
 \qquad
 \dim L_m^{c=-1}=\frac{w_m-\operatorname{tr}(c\mid L_m)}2.
\]
For the free metabelian Lie algebra the elements
$\operatorname{ad}_x^a\operatorname{ad}_y^b[x,y]$, $a+b=m-2$,
form a basis in degree $m\ge2$. Thus the dimension is $m-1$ and the
negative conjugation eigenspace has dimension $\lfloor m/2\rfloor$.
The Hodge argument is unchanged.

As an independent finite check, the program exhausts all binary words
of each length $1$ through $8$, selects the Lyndon words, and counts
them and their signs $(-1)^{\#y}$. It compares these counts with the
Witt formulas and with the explicit metabelian basis. It also verifies
every CM bidegree multiplicity. This is exhaustive at the stated word
lengths, and does not sample Galois cohomology.
\end{proof}

\begin{statement}{F.8}{Proved}{Exact low-level global dimensions and unconditional finiteness}
For either curve, and every finite $S$ with $5\notin S$,
\[
 \dim\operatorname{Sel}_{S,1}=r,
 \qquad \dim\operatorname{Sel}_{S,2}=r+s.
\]
The same identities hold for the metabelian tower. For $E_1$ the first
strict global--local dimension inequality is at class one, for all
$s=0,1,2$. For $E_6$ and $s=0$ it is at class two. These inequalities
prove finiteness of the corresponding integral sets by the unipotent
Albanese method. In particular
$X_1(\mathbb Z[1/S])=\{(-1,0),(0,0),(1,0)\}$ for every $S$.
\end{statement}
\begin{proof}
At every $\ell\ne5$, potential good reduction gives
$H^1(\mathbb Q_\ell,V)=0$: $H^0$ and $H^2$ vanish by weights and local
duality, and the local Euler characteristic is zero. The ordinary
abelian Selmer condition is thus independent of $S$, and F.6 proves
the first identity.

For class two, the involution $[-1]$ fixes $b$ and commutes with Galois.
It acts by $-1$ on $V$ and $+1$ on the central line. Its eigenspaces
therefore give a Galois-equivariant vector-space splitting
\[
 \operatorname{Lie}(U_2)=V\oplus\mathbb Q_5(1),
 \qquad [(v,z),(v',z')]=(0,v\wedge v').
\]
In these coordinates a cocycle consists of $(a,b)$ with
$da=0$ and $db=-\tfrac12a\mathbin{\cup}a$.
For rank zero only $a=0$ occurs. For rank one take the cocycle of the
rational point $P=(12,36)$ and scale its coordinates by
$(a,b)\mapsto(ta,t^2b)$. This supplies an algebraic lift for every
$t\in\mathbb Q_5$ in the full abelian Selmer line; it remains
crystalline at $5$, since $P$ is integral there and the scaling is an
endomorphism of the class-two Galois group object. Consequently the
global class-two fibres before imposing the remaining local conditions
are torsors under
\[
 H^1(G_{\mathbb Q,T},\mathbb Q_5(1))
   =\mathbb Z[1/T]^\times\otimes\mathbb Q_5.
\]
There is no stabilizer quotient: the relevant twisted $H^0$ groups
vanish, since their central graded pieces have negative weights.

The valuation map from this unit space to
$\bigoplus_{\ell\in T}\mathbb Q_5$ is an isomorphism. Crystallinity
at $5$ fixes the $5$-valuation. At $\ell\in B_d\setminus S$, the
integral local Kummer images are finite, and fixing a reduction type
fixes the corresponding central valuation. At a prime in $S$ the
cuspidal component has Zariski closure the full affine central line:
the puncture is rational, and its leading commutator is the nonzero
generator of $\mathbb Q_5(1)$, so its unbounded integral contact orders
are Zariski dense in that line. Thus the type which is cuspidal at
every prime in $S$ leaves exactly $s$ unit parameters. Types which
are noncuspidal at some of them leave fewer parameters. The unit
valuation isomorphism both satisfies the required constants at bad
primes and proves that the maximal dimension is exactly $r+s$.
The finiteness of the non-$5$ integral local images is
\citet[Theorem 0.1 and Corollary 0.2]{FE_KimTamagawa2008}.

F.7 gives the local dimensions $1,2$ at these levels. A proper
algebraic closed subset of the local period space is cut out by a
nonzero regular function. The pullback is a Coleman function; the
unipotent Albanese map is Zariski dense on each residue disk, so it
does not vanish identically there. On the compact integral disks a
nonzero such analytic function has finitely many zeroes. There are
finitely many disks and finitely many reduction types. This proves
finiteness. The original density theorem is
\citet[Theorem 1, p.~93; finiteness deduction in \S3]{FE_Kim2009}.
His indexing starts with $U_2$ abelian, so his $U_{n+1}$ is our $U_n$.
The elliptic rank-at-most-one application appears explicitly in
\citet[Corollary 0.3 and its proof in \S2]{FE_KimTamagawa2008}, whose
$U_3$ is our class two. The modern refined formulation also follows
from \citet[Theorem A, Remark 1.4, and Lemma 3.2]{FE_LLM2023}.

For rank one the abelian localization is nonzero: a non-torsion
rational point has nonzero elliptic $5$-adic logarithm. Thus class one
does not give a dimension inequality. At $s=0$ class two does. The
description of $E_1(\mathbb Q)$ in F.6 gives the last assertion.
\end{proof}

\begin{statement}{F.9}{Open}{Higher elliptic dimensions: exact defects and a conditional complete table}
For $n\ge3$ define the actual finite-dimensional arithmetic invariants
\[
 \delta_n(E)=\dim H_f^1\bigl(\mathbb Q,L_n^*(1)\bigr),\qquad
 \delta_n^{\rm met}(E)=
       \dim H_f^1\bigl(\mathbb Q,(L_n^{\rm met})^*(1)\bigr).
\]
The exact strict global graded dimensions are
\[
 \dim H_f^1(\mathbb Q,L_n)=\dim L_n^{c=-1}+\delta_n(E),
 \quad
 \dim H_f^1(\mathbb Q,L_n^{\rm met})
       =\lfloor n/2\rfloor+\delta_n^{\rm met}(E).
\]
Consequently for $n\ge2$ there are the unconditional bounds
\[
 \begin{split}
 \dim\operatorname{Sel}_{S,n}
   &\le r+s+\sum_{m=3}^n\bigl(\dim L_m^{c=-1}+\delta_m(E)\bigr),\\
 \dim\operatorname{Sel}^{\rm met}_{S,n}
   &\le r+s+\sum_{m=3}^n
           \bigl(\lfloor m/2\rfloor+\delta_m^{\rm met}(E)\bigr).
 \end{split}
\]
If the relevant defects vanish through the displayed level, these
bounds are equalities. In particular, subject to that explicitly
stated hypothesis, the exact global dimensions through class eight are
\[
\begin{array}{cc|rrrrrrrr}
 \multicolumn{2}{c|}{\text{full free tower}}&1&2&3&4&5&6&7&8\\
 r&s&&&&&&&&\\\hline
 0&0&0&0&1&3&6&11&20&36\\
 0&1&0&1&2&4&7&12&21&37\\
 0&2&0&2&3&5&8&13&22&38\\
 1&0&1&1&2&4&7&12&21&37\\
 1&1&1&2&3&5&8&13&22&38\\
 1&2&1&3&4&6&9&14&23&39
\end{array}
\]
\[
\begin{array}{cc|rrrrrrrr}
 \multicolumn{2}{c|}{\text{metabelian tower}}&1&2&3&4&5&6&7&8\\
 r&s&&&&&&&&\\\hline
 0&0&0&0&1&3&5&8&11&15\\
 0&1&0&1&2&4&6&9&12&16\\
 0&2&0&2&3&5&7&10&13&17\\
 1&0&1&1&2&4&6&9&12&16\\
 1&1&1&2&3&5&7&10&13&17\\
 1&2&1&3&4&6&8&11&14&18
\end{array}
\]
The first strict dimension inequalities for rank one are at classes
$2,3,4$ for $s=0,1,2$. The class-two conclusion is unconditional;
the class-three conclusion uses $\delta_3(E_6)=0$, and the class-four
conclusion uses $\delta_3(E_6)=\delta_4(E_6)=0$. The full and metabelian
groups agree through class four. This audit does not establish those
vanishings, or the defects through class eight; the corresponding
unconditional numerical higher Selmer dimensions remain unresolved here.
\end{statement}
\begin{proof}[Proof of the reductions and conditional assertions]
For a geometric representation $W$, Poitou--Tate duality gives
\[
 h_f^1(\mathbb Q,W)=h^0(\mathbb Q,W)
 +h_f^1(\mathbb Q,W^*(1))-h^0(\mathbb Q,W^*(1))
 +h_f^1(\mathbb Q_5,W)-\dim W^{c=1}.
\]
This precise formula is recorded in
\citet[(5.6)]{FE_LLM2023}. For $W=L_n$, $n\ge3$, both global
$H^0$ terms vanish by weights, and $h_f^1(\mathbb Q_5,W)=\dim W$
by F.7. This proves the exact graded formula, including its defect.
For $L_2=\mathbb Q_5(1)$ the dual invariant term is one and the strict
global group is zero; it cannot be treated by the $n\ge3$ formula.
Allowing denominator primes supplies the class-two unit term $s$ of F.8.

Here potential good reduction is useful beyond dimension estimates.
For $n\ge3$ and $\ell\ne5$,
\[
 H^i(\mathbb Q_\ell,L_n)=0\quad(i=0,1,2).
\]
Indeed after a finite extension the representations are unramified and
pure of weight $-n$; the dual twist has weight $n-2$. Neither weight
is zero. Thus the two invariant groups vanish, and local duality and
Euler characteristic give the assertion. The same argument works for
the metabelian layers. At $5$, $F^0L_n=0$ and the dual twist has
$F^0$ equal to its whole de Rham space, with no invariant vectors.
It follows that
\[
 H^1(\mathbb Q_5,L_n)=H_f^1(\mathbb Q_5,L_n),\qquad
 H_f^1(\mathbb Q_5,L_n^*(1))=0.
\]
Poitou--Tate therefore identifies
\[
 H^2(G_{\mathbb Q,T},L_n)^*
       =H_f^1(\mathbb Q,L_n^*(1)).
\]
There is no dependence on the extra finite primes of $T$ in these
groups: the displayed local groups at such primes vanish.

In a central extension
$1\to L_n\to U_n\to U_{n-1}\to1$, the obstruction to lifting a
Galois cocycle lies in this $H^2$, and every nonempty fibre is a
torsor under $H^1(G_{\mathbb Q,T},L_n)$, since the twisted $H^0$
stabilizers vanish. This proves the unconditional upper bound by
induction. If $\delta_n=0$, all lower global classes lift. At a
non-$5$ place their lift is unique. At $5$ the crystalline condition
also lifts: the de Rham quotients
$U_n^{\rm dR}/F^0\to U_{n-1}^{\rm dR}/F^0$ are surjective by
strictness of the Hodge filtration, and the whole central $H^1$ is
crystalline. Thus it is the fibre of the local lift.
Consequently no additional local equation appears in degree $n$.
The map of refined Selmer schemes is then a torsor, locally on the
base, under the vector group of dimension $\dim L_n^{c=-1}$.
Starting with the maximal class-two component of dimension $r+s$
from F.8 proves equality, rather than merely a conditional upper
bound. The identical argument proves the metabelian assertion.
Summing the rows of F.7 gives the two tables.

For completeness the remaining groups have an explicit CM description.
Over $K=\mathbb Q(i)$ write $V=\psi\oplus\overline\psi$ and
$\psi\overline\psi=\chi_5$. The multiplicity of
$\psi^a\overline\psi^b$ in $L_{a+b}$ is
\[
 \ell_{a,b}=\frac1{a+b}
   \sum_{e\mid\gcd(a,b)}\mu(e)
                  \binom{(a+b)/e}{a/e},\qquad a,b\ge1.
\]
It is one in the metabelian layer. The dual-twist characters are
$\psi^{1-a}\overline\psi^{1-b}$, with these multiplicities;
restriction to $K$ and the order-two descent involution recover the
groups over $\mathbb Q$. The program lists every pair $(a,b)$ through
$a+b=8$. This reduces the missing arithmetic to a finite, named list
of Hecke-character Selmer groups, but Hodge weights alone do not
force their vanishing. In particular $L_3=V(1)$, so already
$\delta_3(E)=h_f^1(\mathbb Q,V^*)$ is a distinct problem from
the Mordell--Weil rank calculation.

There are meaningful alternative theorems, but they do not certify the
uncomputed rows. \citet[Definition 2.2 and Corollary 2.9]{FE_Bars2007}
proves the requisite second-cohomology vanishing under geometric
regularity of $5$ for $E$ and $K(E[5])$; that regularity has not been
verified here. \citet[Theorem 0.1]{FE_Kim2010CM} proves eventual
global--local dimension inequality unconditionally for a smaller CM
quotient $W$; his Theorem 0.2 gives an explicit bound subject to
specified nonvanishing values of $p$-adic $L$-functions. The quotient
$W$ kills the Lie bidegrees with both entries at least two and is a
quotient of the metabelian group. Thus it gives an alternative
unconditional eventual finiteness argument for both curves, after
enlarging $S$ to include bad primes if necessary. Its existence theorem
does not supply a first level at most eight or the full-table defects.

Any conditional strict inequality established above implies finiteness
by the same density and analytic argument as F.8; the general effective
refinement is \citet[Theorem 6.2.1(A),(B), and \S6.3]{FE_Betts2023}.
For transfer to LM6, the lemma must retain the specified local
conditions, the rank or finite-$\Sha$ input, the actual Galois action,
and every needed $\delta_m=0$ hypothesis. Replacing these hypotheses
by an assertion about a normalizer Lie algebra supplies no such
arithmetic vanishing theorem.
\end{proof}
